% ================================================================
%  conteudo/revisao.tex
% ================================================================

\section{Revisão da Literatura}

\subsection{Leveduras}

\subsubsection{\textit{Saccharomyces cerevisiae}}

\textit{Saccharomyces cerevisiae} é uma levedura eucariota pertencente ao filo Ascomycota, amplamente empregada em processos biotecnológicos industriais. Sua capacidade de fermentar açúcares como glicose e sacarose, produzindo etanol e dióxido de carbono, a torna o microrganismo modelo por excelência na indústria de bebidas, panificação e biocombustíveis (WALKER, 1998). A linhagem ATCC~7754 é uma cepa de referência laboratorial amplamente utilizada em estudos de fisiologia e controle de bioprocessos.

Do ponto de vista fisiológico, \textit{S. cerevisiae} é um organismo facultativo, capaz de crescer tanto na presença quanto na ausência de oxigênio. Em condições aeróbias e com baixa disponibilidade de glicose, o metabolismo é predominantemente respiratório, com maior produção de biomassa. Já em condições de excesso de glicose ou anaerobiose, o efeito Crabtree direciona o metabolismo para a fermentação alcoólica, mesmo na presença de oxigênio (BASSO et al., 2011; GOMBERT; MELO, 2000).

\subsection{Métodos de Quantificação de Biomassa}

\subsubsection{Densidade Óptica}

A turbidimetria é um método indireto de quantificação de biomassa baseado na capacidade das células microbianas de dispersar a luz incidente em uma suspensão. A absorbância medida a 600~nm é proporcional à concentração de células na amostra, desde que se trabalhe dentro da faixa linear do espectrofotômetro, geralmente entre 0,1 e 0,8 UA (GLAZYRINA et al., 2010). Para amostras com absorbância superior a esse intervalo, é necessário realizar diluições da suspensão a fim de garantir a linearidade da leitura.

A relação entre a absorbância e a concentração celular pode ser expressa pela Lei de Beer-Lambert, que estabelece uma proporcionalidade entre a quantidade de luz absorvida e a concentração do soluto. No contexto microbiológico, essa correlação é válida para suspensões homogêneas de células com tamanho e morfologia uniformes, como é o caso das leveduras (REED; NAGODAWITHANA, 1991). A conversão dos valores de transmitância (\%T) em absorbância é realizada conforme a Equação~\ref{eq:absorbancia}.

\begin{equation}
    A = 2 - \log \%T
    \label{eq:absorbancia}
\end{equation}

\subsubsection{Determinação de Biomassa por Massa Seca}

O método gravimétrico da massa seca é considerado o padrão de referência para a quantificação direta de biomassa microbiana. O procedimento consiste na separação das células do meio de cultura por centrifugação, seguida de lavagem com água destilada para remoção de sais e componentes do meio, e posterior secagem em estufa até massa constante (SCHMIDELL et al., 2001). A concentração de biomassa, expressa em g/L, é calculada conforme as Equações~\ref{eq:massa_seca} e \ref{eq:concentracao}.

\begin{equation}
    m_{\text{seca}} = m_{\text{total}} - m_{\text{tubo}}
    \label{eq:massa_seca}
\end{equation}

\begin{equation}
    \text{Concentração} = \frac{m_{\text{seca}}}{V}
    \label{eq:concentracao}
\end{equation}

Apesar de ser um método preciso, a determinação da massa seca é demorada e consome volume de amostra, o que limita sua aplicação em acompanhamentos cinéticos de alta frequência. Por isso, a curva padrão que correlaciona massa seca e absorbância é uma ferramenta essencial em laboratórios de bioprocessos (GOMBERT; MELO, 2000).
